We start our discussion from Lemma~\eqref{lem:central-newtons-formula-for-powers},
which yields backward Newton's interpolation formula for powers.
For example,
\begin{example}
  For integer $0 \le t \le 3$, we have
  \begin{align*}
    t=0:\quad n^3
    &=
    \bigl(1+\tbinom{n}{-1}\bigr)\cdot0
    +\tfrac{1}{2}(1+2n)\cdot\tfrac{1}{4}
    +\tfrac{1}{2}(2+n+n^2)\cdot0 \\
    &+\tfrac{1}{48}(21+46n+12n^2+8n^3)\cdot6,
    \\
    t=1:\quad n^3
    &=
    \bigl(1+\tbinom{n-1}{-1}\bigr)\cdot1
    +\tfrac{1}{2}(-1+2n)\cdot\tfrac{13}{4}
    +\tfrac{1}{2}(2-n+n^2)\cdot6 \\
    &+\tfrac{1}{48}(-21+46n-12n^2+8n^3)\cdot6,
    \\
    t=2:\quad n^3
    &=
    \bigl(1+\tbinom{n-2}{-1}\bigr)\cdot8
    +\tfrac{1}{2}(-3+2n)\cdot\tfrac{49}{4}
    +\tfrac{1}{2}(4-3n+n^2)\cdot12 \\
    &+\tfrac{1}{48}(-87+94n-36n^2+8n^3)\cdot6,
    \\
    t=3:\quad n^3
    &=
    \bigl(1+\tbinom{n-3}{-1}\bigr)\cdot27
    +\tfrac{1}{2}(-5+2n)\cdot\tfrac{109}{4}
    +\tfrac{1}{2}(8-5n+n^2)\cdot18 \\
    &+\tfrac{1}{48}(-225+190n-60n^2+8n^3)\cdot6.
  \end{align*}
\end{example}

Therefore, by applying
Lemma~\eqref{lem:central-factorials-binomial-form}
to Newton's interpolation formula~\eqref{lem:central-newtons-formula-for-powers},
for non-negative integers $n,m$ and an arbitrary integer $t$ yields
\begin{align*}
  n^m
  &= \frac{\centralFactorial{(n-t)}{0}}{0!} \delta^{0} t^m
  + \sum_{k=1}^{m} \frac{n-t}{k} \binom{n+t+\frac{k}{2}-1}{k-1} \delta^{k} t^m \\
  &= t^m + \sum_{k=1}^{m} (n-t) \binom{n-t+\frac{k}{2}-1}{k-1} \frac{\delta^{k} t^m}{k}.
\end{align*}
Note that we isolate the term $\frac{\centralFactorial{(n-t)}{0}}{0!} \delta^{0} t^m$
to avoid division by zero.
By expanding the brackets, we have
\begin{align*}
  n^m
  = t^m + \sum_{k=1}^{m} \frac{\delta^{k} t^m}{k}
  \left[
    n\binom{n-t+\frac{k}{2}-1}{k-1}
    -
    t\binom{n-t+\frac{k}{2}-1}{k-1}
  \right].
\end{align*}
Now we notice that sum above
has an additional factor $n$ in its binomial basis
$n\binom{n-t+\frac{k}{2}-1}{k-1}$,
preventing us from using centered
hockey-stick identity~\eqref{lem:centered-hockey-stick-identity}
for closed form.
Thus, consider the following Lemma
\begin{lemma}[Binomial decomposition]
  \label{lem:binomial-decomposition}
  For non-negative integers $r,n,m$
  \begin{align*}
    n \binom{n+r}{m} = (m+1) \binom{n+r}{m+1} - (r-m) \binom{n+r}{m}.
  \end{align*}
  \begin{proof}
    By expanding the brackets yields,
    \begin{align*}
      n \binom{n+r}{m}
      = m \binom{n+r}{m+1} + \binom{n+r}{m+1} - r \binom{n+r}{m}
      + m \binom{n+r}{m}.
    \end{align*}
    Recall the extraction property of binomial coefficients
    $\binom{n}{k+1} = \frac{n-k}{k+1} \binom{n}{k}$.
    Therefore, by extraction property
    \begin{align*}
      \binom{n+r}{m+1} = \frac{n+r-m}{m+1} \binom{n+r}{m}.
    \end{align*}
    Thus,
    \begin{align*}
      n \binom{n+r}{m}
      &= m \frac{n+r-m}{m+1} \binom{n+r}{m}
      + \frac{n+r-m}{m+1} \binom{n+r}{m}
      - r \binom{n+r}{m} + m \binom{n+r}{m}.
    \end{align*}
    By factoring out binomial coefficient $\binom{n+r}{m}$ yields
    \begin{align*}
      n \binom{n+r}{m}
      &= \binom{n+r}{m} \left[ m \frac{n+r-m}{m+1} + \frac{n+r-m}{m+1} - r + m \right] \\
      &= \binom{n+r}{m} \left[ (m+1) \frac{n+r-m}{m+1} - r + m \right] \\
      &= \binom{n+r}{m}n.
    \end{align*}
    This completes the proof.
  \end{proof}
\end{lemma}
Hence, by setting
$r \rightarrow -t+\frac{k}{2} -1$,
and $m \rightarrow k-1$
into Lemma~\eqref{lem:binomial-decomposition} gives
\begin{align*}
  n \binom{n-t+\frac{k}{2} -1}{k-1}
  = k \binom{n-t+\frac{k}{2}-1}{k}
  + \left [t+\frac{k}{2} \right ]
  \binom{n-t+\frac{k}{2}-1}{k-1}.
\end{align*}
Because, by binomial decomposition~\eqref{lem:binomial-decomposition}, we have
\begin{align*}
  n \tbinom{n-t+\frac{k}{2} -1}{k-1}
  &=
  k \tbinom{n-t+\frac{k}{2} -1}{k-1+1}
  - \left [-t+\tfrac{k}{2} -1-(k-1)\right ]
  \tbinom{n-t+\frac{k}{2} -1}{k-1} \\
  &= k \tbinom{n-t+\frac{k}{2}-1}{k}
  - \left [-t-\tfrac{k}{2} \right ]
  \tbinom{n-t+\tfrac{k}{2}-1}{k-1} \\
  &= k \tbinom{n-t+\tfrac{k}{2}-1}{k}
  + \left [t+\tfrac{k}{2} \right ]
  \tbinom{n-t+\tfrac{k}{2}-1}{k-1}.
\end{align*}
Thus,
\begin{align*}
  n^m
  = t^m
  + \sum_{k=1}^{m} \frac{\delta^{k} t^m}{k}
  \left[
    k \binom{n-t+\frac{k}{2}-1}{k}
    +
    \left[t+\frac{k}{2} \right]
    \binom{n-t+\frac{k}{2}-1}{k-1}
    -
    t \binom{n-t+\frac{k}{2}-1}{k-1}
  \right].
\end{align*}
By collapsing the terms $t \binom{n-t+\frac{k}{2}-1}{k-1}$, and by simplifying
the factors $k$ yields
\begin{align*}
  n^m
  = t^m
  + \sum_{k=1}^{m}
  \left[
    \binom{n-t+\frac{k}{2}-1}{k}
    +
    \frac{1}{2}
    \binom{n-t+\frac{k}{2}-1}{k-1}
  \right] \delta^{k} t^m.
\end{align*}
We move the term $t^m$ under the summation, thus
\begin{align*}
  n^m
  = \sum_{k=0}^{m}
  \left[
    \binom{n-t+\frac{k}{2}-1}{k}
    +
    \frac{1}{2}
    \binom{n-t+\frac{k}{2}-1}{k-1}
  \right] \delta^{k} t^m.
\end{align*}
Now, we simplify the binomial basis
$\binom{n-t+\frac{k}{2}-1}{k}+\frac{1}{2}\binom{n-t+\frac{k}{2}-1}{k-1}$
to get rid of fractions inside of it.
Let $a=n-t+\frac{k}{2}-1$, then
\begin{align*}
  \binom{a}{k} + \frac{1}{2} \binom{a}{k-1}
  = \frac{1}{2}
  \left[
    2 \binom{a}{k} + \binom{a}{k-1}
  \right]
  = \frac{1}{2}
  \left[
    \binom{a}{k} + \binom{a+1}{k}
  \right].
\end{align*}
Therefore,
\begin{align*}
  n^m
  = \frac{1}{2}
  \sum_{k=0}^{m}
  \left[
    \binom{n-t+\frac{k}{2}-1}{k}
    +
    \binom{n-t+\frac{k}{2}}{k-1}
  \right] \delta^{k} t^m.
\end{align*}
For example,
